\documentclass[12pt]{article}
\usepackage{amsfonts,amsthm,amsmath,amssymb,mathtools}
\usepackage{mathrsfs}
\usepackage{enumitem}
\usepackage{tikz-cd}
\usepackage{geometry}
\usepackage{mlmodern}

\geometry{letterpaper, portrait, margin=1in}

\setcounter{section}{0}

\renewcommand{\thesection}{\S\arabic{section}}
\renewcommand{\thesubsection}{\arabic{section}}

\newtheorem{thm}{Theorem}
\newenvironment{theorem}[1]{
	\renewcommand{\thethm}{#1}
	\begin{thm}
		}{\end{thm}}

\newtheorem{lma}{Lemma}
\newenvironment{lemma}[1]{
	\renewcommand{\thelma}{#1}
	\begin{lma}
		}{\end{lma}}

\theoremstyle{definition}
\newtheorem{ex}{}
\newenvironment{exercise}[1]{
	\renewcommand{\theex}{#1}
	\begin{ex}
		}{\end{ex}}

\title{Homework 1}
\author{Connor Mooneyhan}
\date{September 4, 2025}

\begin{document}

\maketitle

\begin{ex}
	Show that there are no integers $x$, $y$, and $z$ so that \begin{equation*}
		x^3 + y^3 + z^3 = 4.
	\end{equation*}
\end{ex}
\begin{proof}
	First, observe the following cubes considered mod 9: \begin{align*}
		0^3 & \equiv 0 \pmod 9  \\
		1^3 & \equiv 1 \pmod 9  \\
		2^3 & \equiv 8 \pmod 9  \\
		3^3 & \equiv 0 \pmod 9  \\
		4^3 & \equiv 1 \pmod 9  \\
		5^3 & \equiv 8 \pmod 9  \\
		6^3 & \equiv 0 \pmod 9  \\
		7^3 & \equiv 1 \pmod 9  \\
		8^3 & \equiv 8 \pmod 9.
	\end{align*}
	Thus the only values congruent to a cube mod 9 are 0, 1, and 8.

	Now consider the equation mod 9.
	We would need to have $a,b,c \in \left\lbrace 0, 1, 8 \right\rbrace$ such that $a + b + c \equiv 4 \bmod 9$, which is impossible since \begin{align*}
		0 + 0 + 0 & \equiv 0 \pmod 9  \\
		0 + 1 + 1 & \equiv 2 \pmod 9  \\
		0 + 8 + 8 & \equiv 7 \pmod 9  \\
		0 + 0 + 1 & \equiv 1 \pmod 9  \\
		0 + 0 + 8 & \equiv 8 \pmod 9  \\
		0 + 1 + 8 & \equiv 0 \pmod 9  \\
		1 + 1 + 1 & \equiv 3 \pmod 9  \\
		1 + 1 + 8 & \equiv 1 \pmod 9  \\
		1 + 8 + 8 & \equiv 8 \pmod 9  \\
		8 + 8 + 8 & \equiv 6 \pmod 9.
	\end{align*}
\end{proof}

\begin{ex}\
	\begin{enumerate}[label=(\alph*)]
		\item Show that the only integer solutions to $w^2 = x^2 - x + 1$ are those with $x = 0$ and $x = 1$. (Hint: Write the equation as $4w^2 - (4x^2 - 4x + 1) = 3$.)
		\item Find, with proof, all integer solutions to the equation $y^2 = x^3 - x^2 + x$.
	\end{enumerate}
\end{ex}
\begin{proof}
	\begin{enumerate}[label=(\alph*)]
		\item Multiply the equation by 4 to get \begin{align*}
			               & 4w^2 = 4x^2 - 4x + 4       \\
			      \implies & 4w^2 - 4x^2 + 4x - 1 = 3   \\
			      \implies & 4w^2 - (2x-1)^2 = 3        \\
			      \implies & (2w-(2x-1))(2w+(2x-1)) = 3 \\
			      \implies & (2w-2x+1)(2w+2x-1) = 3.
		      \end{align*}
		      Since 3 is prime, this means that the factors are either 1 and 3 or $-1$ and $-3$.
		      First, we evaluate (by subtracting equations) \begin{align*}
			               & 2w - 2x + 1 = 1 \\
			               & 2w + 2x - 1 = 3 \\
			      \implies & 4x - 2 = 2      \\
			      \implies & x = 1,
		      \end{align*}
		      then \begin{align*}
			               & 2w - 2x + 1 = 3 \\
			               & 2w + 2x - 1 = 1 \\
			      \implies & 4x - 2 = -2     \\
			      \implies & x = 0.
		      \end{align*}
		      For the negative variants, we have \begin{align*}
			               & 2w - 2x + 1 = -1 \\
			               & 2w + 2x - 1 = -3 \\
			      \implies & 4x - 2 = -2      \\
			      \implies & x = 0
		      \end{align*}
		      and \begin{align*}
			               & 2w - 2x + 1 = -3 \\
			               & 2w + 2x - 1 = -1 \\
			      \implies & 4x - 2 = 2       \\
			      \implies & x = 1,
		      \end{align*}
		      so the only integer solutions have $x=0$ or $x=1$.
		\item We will denote a solution pair where $x = x_1$ and $y = y_1$ as $(x_1, y_1)$.
		      By factoring, we get $y^2 = x(x^2 - x + 1)$.
          Immediately, we see that if $x = 0$, $(0, 0)$ is the only corresponding solution, and if $x = 1$, $(1, 1)$ and $(1, -1)$ are the corresponding solutions.
		      We can reject any cases where $x$ is negative, since $y^2$ is positive and $x^2 - x + 1$ is always positive since its discriminant is negative and it has a positive value at $x = 0$, so $x$ must be positive as well to yield the positive $y^2$.
		      Note that $\gcd(x, x^2 - x + 1) = 1$ by Bezout's lemma since $(1 - x)\cdot x + 1 \cdot (x^2 - x + 1) = 1$.
		      Then $y^2$ is the product of two relatively prime integers.

		      Let $\prod p_i^{k_i}$ and $\prod p_j^{\ell_j}$ be the prime factorizations of $x$ and $x^2 - x + 1$, respectively.
		      Since there are no shared primes between the two, all $k_i$ and $\ell_j$ are even, and thus $x = a^2$ and $x^2-x+1 = b^2$ for some integers $a$ and $b$.
		      We proceed to show that $x^2 - x + 1$ is never a square when $x > 1$.

		      Indeed, \begin{align*}
			               & 1 < x              \\
			      \implies & -x + 1 < 0         \\
			      \implies & x^2 - x + 1 < x^2,
		      \end{align*}
		      and \begin{align*}
			      (x-1)^2 & = x^2 - 2x + 1   \\
			              & < x^2 - x + 1,
		      \end{align*}
          so \begin{equation*}
            (x-1)^2 < x^2 - x + 1 < x^2.
          \end{equation*}
          Since $f: x \mapsto x^2$ is an increasing function over the nonnegative integers, we can extend this to \begin{equation*}
            m^2 \leq (x-1)^2 < x^2 - x + 1 < x^2 \leq n^2
          \end{equation*}
          where $0 \leq m \leq x-1$ and $x \leq n$.
          This tells us that given any integer $y$ (since $(-y)^2 = y^2$), either $y^2 < x^2 - x + 1$ or $x^2 - x + 1 < y^2$, and thus $x^2 - x + 1$ is not the square of an integer when $x > 1$.

          We can now comfortably say, then, that the only solutions occur when $0 \leq x \leq 1$, namely the three solutions $(0, 0)$, $(1, 1)$, and $(1, -1)$ that we found at the beginning.
	\end{enumerate}

\end{proof}

\begin{ex}
	$\ast$\ Prove that there are infinitely many integer solutions to $x^2 - 5y^2 = 1$.
	Here are the first few solutions:
	\begin{equation*}
		\begin{array}{cc}
			x      & y          \\
			\hline
			1      & 0          \\
			9      & \pm 4      \\
			161    & \pm 72     \\
			2889   & \pm 1292   \\
			51841  & \pm 23184  \\
			930249 & \pm 416020 \\
		\end{array}
	\end{equation*}
	Based on these examples, try to figure out how to generate the next solution from the previous one, and use induction to prove that there are infinitely many.
\end{ex}
\begin{proof}
	We claim that the following relationship holds for the above pairs of solutions: \begin{align}
		x_i & = 9x_{i-1} + 20y_{i - 1}\label{3.1} \\
		y_i & = 4x_{i-1} + 9y_{i - 1}.\label{3.2}
	\end{align}
	The induction proof works whether or not this relationship holds for the above, so we will not prove the relationship here and will instead skip to the induction proof.

	Assume that there exists some pair of solutions $(x_{i-1}, y_{i-1})$, and let $(x_i, y_i)$ be determined by $\eqref{3.1}$ and \eqref{3.2}.
	Then \begin{align*}
		x_i^2 - 5y_i^2 & = (9x_{i - 1} + 20y_{i - 1})^2 - 5(4x_{i-1} + 9y_{i-1})^2                                          \\
		               & = 81x_{i-1}^2 + 360x_{i-1}y_{i-1} + 400y_{i-1}^2 - 5(16x_{i-1}^2 + 72x_{i-1}y_{i-1} + 81y_{i-1}^2) \\
		               & = 81x_{i-1}^2 + 360x_{i-1}y_{i-1} + 400y_{i-1}^2 - 80x_{i-1}^2 - 360x_{i-1}y_{i-1} - 405y_{i-1}^2  \\
		               & = x_{i-1}^2 - 5y_{i-1}^2                                                                           \\
		               & = 1.
	\end{align*}
	Certainly, $x_i$ and $y_i$ are integers since they are $\mathbb{Z}$-linear combinations of integers.
	By induction, then, taking $(x_1, y_1)$ to be $(1, 0)$, there is an integer solution $(x_i, y_i)$ for each $i \in \mathbb{Z}^+$, and thus there are infinitely many.
\end{proof}

\end{document}
